\documentclass[11pt]{article}
\usepackage[utf8]{inputenc}
\usepackage{amsmath, amssymb}
\usepackage{graphicx}
\usepackage{natbib}
\usepackage{geometry}
\usepackage{setspace}
\geometry{margin=1in}
\usepackage{hyperref}
\usepackage{titlesec}
\usepackage{booktabs}
\usepackage{caption}
\usepackage{float}
\usepackage{threeparttable}
\usepackage[english]{babel}


\title{UEFA vs FIFA Tie Break: Which Creates More Suspense?}
\author{Alessandro Di Mattia and Alex Krumer}

\date{}

\begin{document}

\maketitle

\begin{center}
\includegraphics[width=0.3\textwidth]{images/logo.png}

\vspace{0.5em}

This version: November 2026 \\
Date this version was printed: \today
\end{center}

\begin{abstract}
This paper studies how tie-breaking rules affect group-level suspense in international football tournaments. We compare the goal-difference criterion used by FIFA in the World Cup with the head-to-head criterion used by UEFA in the European Championship. For that, we construct a binary suspense indicator that equals one when a single goal could change the set of teams qualifying for the knockout stage. Applying both rules to the same groups, we find that FIFA rules generate significantly more group-level suspense in the World Cup ($p < 0.01$), while in the European Championship the difference is not significant. A logit regression across the two tournaments that controls for competitive balance shows that World Cup goals are less likely to change the set of qualifying teams. This is consistent with the FIFA criterion spreading uncertainty more evenly over the group stage instead of concentrating it in a few decisive goals. We also show that under FIFA rules suspense remains higher throughout the 90 minutes of play. Overall, our results suggest that the goal-difference criterion keeps qualification open for longer in the World Cup.
\end{abstract}

\noindent \textbf{Keywords:} Consumption of sport, suspense, OR in sports, group stage, football

\vspace{1em}

\noindent \textbf{Address for correspondence:}

\noindent Alessandro Di Mattia, Faculty of Business Administration and Social Sciences, Molde University College, Britvegen 2, Molde, 6402, Norway, alessandro1.e.g.d.di.mattia@himolde.no.

\noindent Alex Krumer (corresponding author), Faculty of Business Administration and Social Sciences, Molde University College, Britvegen 2, Molde, 6402, Norway, alex.krumer@himolde.no.

\newpage

\section{Introduction}

Retaining the attention of spectators has become both harder and more important in modern sports, and in football in particular. Fans today can choose among many forms of entertainment, so keeping them engaged is key for the success and profitability of sport organizations. Apart from loyalty to a specific team, fans enjoy \emph{non-instrumental information}, such as the suspense and uncertainty of a match, which satisfies basic psychological needs \citep{ely2015suspense, zillmann2013psychology}. This value often drives interest regardless of the final result, since it comes from the events that lead to that result \citep{vorderer2004enjoyment}.

The economic importance of fan engagement is reflected in the audiences and revenues of major tournaments. For instance, the 2022 FIFA World Cup in Qatar was watched for at least one minute by approximately 2.87 billion viewers, while more than 5 billion people followed it across broadcast, digital, and social media platforms, including 2.9 billion via linear TV and 2.7 billion through digital streaming \citep{fifa2022engagement, fifa2022audience}. Broadcasting and commercial rights from this World Cup cycle generated more than USD 6.3 billion for FIFA, which is 83\% of the organization's revenue \citep{fifa2022finance}.

The UEFA European Championship also attracts large audiences. Euro 2024 had a cumulative stadium attendance of 2.67 million across its 51 matches and 5.7 billion television viewers worldwide \citep{uefa2024numbers, uefa2024remember}. These figures show that audience attention translates directly into revenue.

Given these stakes, sports organizers should not only increase the number of matches but also design competitions that keep suspense and engagement high. Low-stakes games may reduce interest, while rules that preserve uncertainty, such as tie-breaking rules, can keep fans interested throughout the tournament. Therefore, the way tie-breaking procedures affect suspense matters not only for sporting fairness but also financially, since it affects consumer demand, broadcaster revenues, and sponsor value.

In this paper, we study how different tie-breaking rules affect suspense in international football tournaments. More specifically, we compare the two main approaches used by governing bodies: FIFA's use of overall goal difference and goals scored, and UEFA's use of head-to-head results. We introduce a binary suspense indicator at the group level that captures whether a single goal could change the set of teams qualifying for the knockout stage. To the best of our knowledge, this measure is new in the sports economics literature. Our main result is that FIFA's goal-difference criterion generates significantly more group-level suspense in the World Cup. We also find that under FIFA rules suspense remains higher throughout the 90 minutes of play. The same advantage does not appear in the European Championship, where the choice of rule does not significantly affect suspense. This suggests that the effect is specific to the World Cup.

\vspace{0.5em}

The remainder of the paper is organized as follows. Section~\ref{sec:literature} reviews the related literature. Section~\ref{sec:data} presents the data, and Section~\ref{sec:estimation} describes the methodological framework. Section~\ref{sec:results} reports the results, and Section~\ref{sec:conclusion} concludes. Additional analyses, including correlation tables and suspense trajectories by group type, are reported in the Appendix.

\section{Literature Review}
\label{sec:literature}

Our paper is related to the literature on tournament design and on the drivers of sport consumption. Recent studies in sports economics have pointed out several flaws in the rules and formats of international football competitions and have shown how these flaws affect both the fairness of outcomes and the demand for the game. For instance, \citet{csato2025increasing} propose a framework to evaluate the cost of stake-less games, in which the result no longer affects qualification, in the new 2026 FIFA World Cup format. \citet{lapre2022quantifying, lapre2023evolution} show that FIFA's seeding system has failed to reduce competitive imbalance, which considerably lowers the chances of some teams to advance. Regarding UEFA regulations, \citet{csato2020neither} shows that flaws in the qualification rules can create situations in which teams have incentives not to win, which reveals incentive compatibility problems in current tournament designs.

To address these weaknesses, researchers have proposed several solutions, from probabilistic models that improve the predictive accuracy of ranking algorithms \citep{szczecinski2022fifa} to optimal schedules that reduce collusion and stake-less games \citep{chater2021fixing}. Other studies redesign tournament formats and rules to remove misaligned incentives, for example through alternative qualification systems \citep{guajardo2024tournament}, dynamic three-team formats \citep{troyano2023new}, and double-elimination structures \citep{renno2023double}.

\vspace{0.5em}

A related issue is the choice of tie-breaking rules, which directly affects the fairness and the stakes of group-stage matches. For example, \citet{berker2014tie} shows that under certain rules, the ranking of a team can depend on the outcome of a match in which it does not play. \citet{csato2023avoid} shows that the probability of uncompetitive matches, where teams have no incentive to exert full effort, depends strongly on the tie-breaking criteria. Building on this, \citet{csato2026match} show that replacing goal difference with head-to-head results as the primary tie-breaking criterion in the 2026 FIFA World Cup increases the share of unimportant last-round matches by approximately 5.6 percentage points. This is because head-to-head rules restrict the set of possible group rankings after the first two rounds and turn offensive-asymmetric matches, in which one team must attack to survive, into unimportant matches in which teams have little incentive to exert full effort. These findings suggest that poorly designed tie-breakers affect not only match outcomes but also the integrity and attractiveness of competitions.

The consequences of such design flaws go beyond sporting fairness and extend to the economics of sport consumption. Formats and rules that create opportunities for collusion or reduce competitiveness lower fan interest, and with it the entertainment value and commercial revenues of tournaments. It is therefore important to understand what drives demand for both live attendance and broadcast audiences. Several studies point to in-game dynamics. \citet{buraimo2020unscripted} show that suspense, surprise, and shock, which capture the tension and emotional intensity of matches, increase demand, particularly at specific moments of play. \citet{alavy2010edge} argue that the progression of the game, and not only its final outcome, matters, with greater uncertainty about the result increasing spectator engagement. \citet{bizzozero2016importance} find a strong link between suspense and surprise on the one hand and demand for entertainment on the other. \citet{flepp2026suspense} show that the relative importance of suspense and surprise depends on the scoring rate of the match. In high-scoring games, surprise dominates because goals often overturn expectations, whereas in low-scoring games sustained suspense is the main driver of engagement. Finally, \citet{buraimo2022armchair} extend this analysis to television audiences and show that demand depends not only on outcome uncertainty but also on the quality of the players and the importance of the match.

Overall, this literature shows that flawed formats and rules raise not only ethical and fairness concerns but also reduce the economic value of competitions by lowering engagement and the attractiveness of matches. We contribute to this literature by examining how institutional features, in particular tie-breaking rules, affect the stakes of a match and its importance in the qualification race. For that, we introduce a group-level dimension of suspense that complements the match-level measures of suspense, surprise, and shock used in the literature. While existing work captures how in-game events change beliefs about the outcome of a single match, our measure captures how the same events change the qualification picture of the whole group, which match-level measures cannot capture.



\section{Data}
\label{sec:data}

Our dataset covers all group-stage matches of the FIFA World Cup and the UEFA European Championship and includes goal-level and minute-by-minute match events. The data were scraped from Wikipedia and cross-checked with official FIFA and UEFA sources. This allowed us to reconstruct how group standings evolved during matches under each tie-breaking criterion and to directly compare the suspense generated by the two systems.

For the FIFA World Cup, we collected data from the 1986 to the 2022 editions. From 1986 to 1994 there were six groups (Groups A to F), and from 1998 onwards eight groups, each with four teams. In the six-group format, the top two teams of each group and the four best third-placed teams advanced to the knockout stage, whereas in the eight-group format only the top two teams of each group qualified. For the final matchday, our dataset includes 381 goals and 6{,}789 minute-by-minute observations.

For the European Championship, we collected data from 1984 to 2024. The first three editions in our sample had two groups (Group 1 and Group 2), the editions from 1996 to 2012 had four groups (Groups A to D), and the editions from 2016 onwards had six groups (Groups A to F). As in the World Cup, in the six-group format the top two teams and the four best third-placed teams qualify for the next round. For the final matchday, we recorded 226 goals and 4{,}044 minute-by-minute observations.

In both competitions, teams received two points for a win until 1992 and three points from 1994 onwards. A draw was worth one point and a loss zero points in all editions. Our analysis focuses on the third and final matchday of the group stage, when the two matches of each group are played simultaneously and qualification is decided.


\subsection{Data Collection}

For every match in both competitions, we collected the following information:

\begin{itemize}
    \item \textbf{Goals scored}: including the scoring team, exact minute, and the evolving scoreline.
    \item \textbf{Team Elo ratings}: Elo ratings as of the day before each tournament began, retrieved from \url{https://www.international-football.net/}.
\end{itemize}

We used the goal data to reconstruct the group standings after each goal under both the FIFA and the UEFA tie-breaking criteria. Elo ratings are used to measure competitive balance at the group and the match level, as described in Section~\ref{sec:results}. The minute-by-minute dataset has one observation per group and minute, from kick-off (minute 0) to minute 90, except for groups in which at least one goal was scored in stoppage time, where observations extend beyond minute 90.


\section{Methodological Framework}
\label{sec:estimation}

Our main contribution is a \textit{binary suspense indicator} at the group level that captures whether a single goal could change the set of teams qualifying for the knockout stage. In other words, a moment is suspenseful when one additional goal in any ongoing match of the group would change which teams advance. We evaluate the indicator in two ways:

\begin{itemize}
    \item \textbf{After each goal scored}, by updating the group standings with the current match results.
    \item \textbf{Minute by minute}, to track the evolution of suspense on a 90-minute scale.
\end{itemize}

This allows us to examine how long a group remains undecided and how often suspense arises under each tie-breaking rule.

For example, in Group F of Euro 2024, Portugal and Turkey were in the qualifying positions before the third and final matchday. Turkey played against the Czech Republic, and the match started 0--0. At that point, a single goal by the Czech Republic would have allowed them to overtake Turkey on points and qualify instead. Since one goal could change which teams qualify, this situation is classified as suspenseful (suspense~$= 1$).

Another example comes from Group A of Euro 2024. On the final matchday, Hungary scored in the 100th minute against Scotland and won 1--0. This goal put Hungary on 3 points with a goal difference of $-3$, which made them one of the best third-placed teams at that moment and therefore a potential qualifier for the knockout stage. Since suspense is evaluated on the basis of the uncertainty during the match, and not on the results of groups played later, this goal is classified as a suspenseful moment.

Group~E of the 2022 FIFA World Cup, with Spain, Germany, Japan, and Costa Rica, illustrates the \textit{qual\_count} variable. Before the final matchday, Spain and Japan held the two qualifying positions, while Germany and Costa Rica were out. Since the two matches, Spain vs.\ Japan and Costa Rica vs.\ Germany, were played simultaneously, the set of qualifying teams changed four times. Germany's goal in minute~10 moved them above Japan into second place. Japan's equaliser against Spain in minute~48 brought them back into a qualifying position and Germany out. Costa Rica's goal in minute~70 put Costa Rica and Japan in the two qualifying positions, and Germany's equaliser in minute~73 changed the set again. In the end, Japan and Spain qualified, but only after \textit{qual\_count} had reached four. This example shows how the FIFA goal-difference criterion can produce repeated and rapid changes in the set of qualifying teams when matches are played simultaneously.

The full logic of the suspense indicator, covering all specific cases under both FIFA and UEFA rules, is described in Appendix~\ref{sec:suspense-criteria}.


\section{Results}
\label{sec:results}

\subsection{Descriptive Statistics}


Table~\ref{tab:desc_stats_goals} reports descriptive statistics for the goal-level dataset and compares World Cup groups, which use FIFA tie-breaking rules, with European Championship groups, which use UEFA rules. The World Cup sample includes 381 goals in 74 groups, and the European Championship sample includes 226 goals in 44 groups. The averages of the main outcome variables show a consistent pattern: \textit{qual.\ changed} is 0.19 in the World Cup and 0.29 in the European Championship, \textit{qual.\ count} is 0.72 and 1.12, and \textit{suspense} is 0.44 and 0.49, respectively. In all three cases, the averages are higher in the European Championship. European Championship teams are also slightly stronger on average, with a mean Elo rating of the top-ranked team in each group of 1998, compared to 1988 in the World Cup.

Table~\ref{tab:desc_stats_mbm} reports the same comparison for the minute-by-minute dataset (6{,}789 World Cup and 4{,}044 European Championship group minutes). The pattern is the same: \textit{qual.\ changed} is 0.15 in the World Cup and 0.24 in the European Championship, \textit{qual.\ count} is 0.48 and 0.69, and \textit{suspense} is 0.50 and 0.57.

\begin{table}[H]
\centering
\caption{Descriptive statistics, goal-level dataset}
\label{tab:desc_stats_goals}
\begin{threeparttable}
\begin{tabular}{lrrrrrrrr}
\toprule
 & \multicolumn{4}{c}{World Cup (FIFA rules)} & \multicolumn{4}{c}{Euro (UEFA rules)} \\
\cmidrule(lr){2-5}\cmidrule(lr){6-9}
 & Mean & SD & Min & Max & Mean & SD & Min & Max \\
\midrule
Goal minute        & 51.99 & 26.89 &  1 &  96 & 52.41 & 26.30 &  2 & 100 \\
Half time          &  1.60 &  0.49 &  1 &   2 &  1.61 &  0.49 &  1 &   2 \\
Qual.\ changed     &  0.19 &  0.39 &  0 &   1 &  0.29 &  0.46 &  0 &   1 \\
Qual.\ count       &  0.72 &  0.99 &  0 &   4 &  1.12 &  1.31 &  0 &   6 \\
Suspense           &  0.44 &  0.50 &  0 &   1 &  0.49 &  0.50 &  0 &   1 \\
Elo (favourite)    & 1904.6 & 106.8 & 1644 & 2169 & 1940.8 &  84.5 & 1671 & 2127 \\
Elo (underdog)     & 1730.4 &  95.3 & 1460 & 1944 & 1812.0 & 102.1 & 1523 & 2040 \\
Elo 1st            & 1987.9 &  77.1 & 1845 & 2169 & 1998.4 &  58.3 & 1850 & 2127 \\
Elo 2nd            & 1855.7 &  63.6 & 1678 & 1986 & 1909.6 &  71.9 & 1766 & 2043 \\
Elo 3rd            & 1777.7 &  61.4 & 1644 & 1920 & 1843.0 &  69.3 & 1671 & 1973 \\
Elo 4th            & 1668.1 &  72.4 & 1460 & 1852 & 1748.6 &  86.0 & 1523 & 1912 \\
\midrule
Observations       & \multicolumn{4}{c}{381} & \multicolumn{4}{c}{226} \\
\bottomrule
\end{tabular}
\begin{tablenotes}
\small
\item \textit{Notes:} Unit of observation is a goal scored on matchday~3. Initial-state rows (minute~$= 0$, one per group) are excluded. \textit{Qual.\ changed} equals~1 when the set of qualifying teams changes from one goal to the next. \textit{Qual.\ count} is the cumulative number of such changes within the group. \textit{Suspense} equals~1 when the next goal would change the qualifying composition. \textit{Elo (favourite)} and \textit{Elo (underdog)} are the Elo ratings of the stronger and weaker team in each match. \textit{Elo 1st}--\textit{Elo 4th} are group-level Elo ranks from highest to lowest.
\end{tablenotes}
\end{threeparttable}
\end{table}

\begin{table}[H]
\centering
\caption{Descriptive statistics, minute-by-minute dataset}
\label{tab:desc_stats_mbm}
\begin{threeparttable}
\begin{tabular}{lrrrrrrrr}
\toprule
 & \multicolumn{4}{c}{World Cup (FIFA rules)} & \multicolumn{4}{c}{Euro (UEFA rules)} \\
\cmidrule(lr){2-5}\cmidrule(lr){6-9}
 & Mean & SD & Min & Max & Mean & SD & Min & Max \\
\midrule
Qual.\ changed  & 0.15 & 0.36 & 0 & 1 & 0.24 & 0.43 & 0 & 1 \\
Qual.\ count    & 0.48 & 0.81 & 0 & 4 & 0.69 & 1.11 & 0 & 6 \\
Suspense        & 0.50 & 0.50 & 0 & 1 & 0.57 & 0.50 & 0 & 1 \\
Elo 1st         & 1984.9 & 74.4 & 1845 & 2169 & 1993.0 & 57.3 & 1850 & 2127 \\
Elo 2nd         & 1862.6 & 62.5 & 1678 & 1986 & 1896.0 & 72.9 & 1766 & 2043 \\
Elo 3rd         & 1780.7 & 64.5 & 1644 & 1920 & 1833.4 & 70.6 & 1671 & 1973 \\
Elo 4th         & 1668.2 & 76.0 & 1460 & 1852 & 1741.4 & 88.5 & 1523 & 1912 \\
\midrule
Observations    & \multicolumn{4}{c}{6{,}789} & \multicolumn{4}{c}{4{,}044} \\
\bottomrule
\end{tabular}
\begin{tablenotes}
\small
\item \textit{Notes:} Unit of observation is a group minute on matchday~3. The dataset spans from minute~0 to minute~90 for each group, extended beyond minute~90 for groups in which at least one goal was scored in extra time. \textit{Qual.\ changed} equals~1 when a goal at that minute changes the set of qualifying teams. \textit{Qual.\ count} is the cumulative number of such changes within the group up to that minute. \textit{Suspense} equals~1 when the next goal would change the qualifying composition. \textit{Elo 1st}--\textit{Elo 4th} are group-level Elo ranks from highest to lowest.
\end{tablenotes}
\end{threeparttable}
\end{table}



\subsection{Statistical Tests}

Table~\ref{tab:tests_diff} reports paired t-tests and Wilcoxon signed-rank tests that compare the group-level mean suspense and mean qualification count under the two rules for all 74 World Cup groups and all 44 European Championship groups. In the World Cup, FIFA rules generate significantly more group-level suspense than UEFA rules ($p<0.01$, both tests), while the number of actual changes in the set of qualifying teams does not differ significantly between the two systems ($p=0.535$). In other words, the FIFA rule increases the number of \textit{potential} changes in qualification, and therefore suspense, without changing the number of \textit{actual} changes.

In the European Championship, neither suspense nor the number of qualification changes differs significantly between the two rules ($p=0.389$ and $p=0.182$, respectively).

\begin{table}[H]
\centering
\small
\caption{Differences between UEFA and FIFA tie-breaking rules: mean differences and statistical tests}
\label{tab:tests_diff}
\begin{threeparttable}
\begin{tabular}{lcc}
\toprule
\multicolumn{3}{c}{\textbf{European Championship}} \\
\midrule
\textbf{Metric} & Mean Diff. & p-value \\
\midrule
Qualification Count & -0.106 & 0.182 (t-test), 0.250 (Wilcoxon) \\
Suspense            & -0.015 & 0.389 (t-test), 0.438 (Wilcoxon) \\
\midrule
\multicolumn{3}{c}{\textbf{World Cup}} \\
\midrule
\textbf{Metric} & Mean Diff. & p-value \\
\midrule
Qualification Count & -0.024 & 0.535 (t-test), 0.605 (Wilcoxon) \\
Suspense            & -0.040 & 0.008 (t-test), 0.010 (Wilcoxon) \\
\bottomrule
\end{tabular}
\begin{tablenotes}[flushleft]
\small
\item \textit{Note:} Reported values are mean differences computed as UEFA $-$ FIFA over all groups
(44 European Championship groups; 74 World Cup groups). The paired t-test evaluates
differences in means under approximate normality; the Wilcoxon signed-rank test provides
a non-parametric check. Zero differences are excluded from the Wilcoxon test.
\end{tablenotes}
\end{threeparttable}
\end{table}




\subsection{Suspense over Time}


Figure~\ref{fig:gls_dist} shows that the distribution of goals over time is broadly similar in the two competitions. In both tournaments, goals are most frequent around the 20th to 30th minute and between the 45th and 55th minute. Differences appear toward the end of matches. In the European Championship there is a clear peak between the 80th and 85th minute, while in the World Cup goals are spread more evenly between the 80th and 90th minute. Overall, World Cup goals are more evenly distributed over the full 100 minutes of play, whereas European Championship goals are more concentrated in specific intervals.

\begin{figure}[H]
    \centering
\includegraphics[width=1\textwidth]{images/gls_dist.png}
    \caption{Distribution of goals by minute in the European Championship and the World Cup.}
    \label{fig:gls_dist}
\end{figure}

Figures~\ref{fig:susp_minute_wc} and~\ref{fig:susp_minute_eu} show the average suspense per minute under each rule for the World Cup and the European Championship, respectively. In the World Cup, FIFA rules generate higher suspense throughout the 90 minutes, and the gap is largest toward the end of the match. In the European Championship, the difference between the two rules is small and, in line with the statistical tests, not significant. This comparison shows that the effect of the tie-breaking rule on suspense is specific to the World Cup and is not a mechanical feature of how the rules work in any tournament.

\begin{figure}[H]
    \centering
    \includegraphics[width=1\textwidth]{images/susp_minute_wc.png}
    \caption{Average suspense per minute under FIFA and UEFA rules, World Cup.}
    \label{fig:susp_minute_wc}
\end{figure}

\begin{figure}[H]
    \centering
    \includegraphics[width=1\textwidth]{images/susp_minute_eu.png}
    \caption{Average suspense per minute under FIFA and UEFA rules, European Championship.}
    \label{fig:susp_minute_eu}
\end{figure}

This pattern is in line with \citet{csato2023avoid}, who uses theoretical arguments and simulations based on the 2022/23 UEFA Nations League to show that ranking goal difference (FIFA) above head-to-head results (UEFA) raises the stakes of matches toward the end of the competition. Figure~\ref{fig:susp_minute_wc} provides empirical support for this mechanism in the World Cup. Suspense trajectories by group type for both tournaments are reported in Appendix~\ref{app:wc_trajectories} and Appendix~\ref{app:euro_trajectories}.





\subsection{Regression Analysis}

We estimate seven models to test whether the FIFA rule is associated with more frequent qualification changes and higher suspense after controlling for competitive balance. For all logit models, we report average marginal effects (AME) with standard errors clustered at the group (year~$\times$~stage) level. Models with match-level Elo controls (Models~1 and 4) use the 607 goals only, since match-level Elo is not defined for the initial state of each group. Models with group-level Elo ranks (Models~2 and 5) also include the 118 initial-state observations, one per group at kick-off, which gives 725 observations.

\paragraph{Model 1: Logit for qualification change, match-level Elo (Table~\ref{tab:reg_grp}).}
The unit of observation is a goal $k$ in group $g$. The dependent variable is \textit{qual\_changed}, a binary indicator equal to~1 when the set of qualifying teams changes from one goal to the next:
\[
\Pr(\textit{qual\_changed}_{gk} = 1)
  = \Lambda\!\left(\alpha + \beta\,\textit{FIFA\_rule}_g
  + \gamma_1\,\textit{elo\_fav}_{gk} + \gamma_2\,\textit{elo\_und}_{gk}\right)
\]
where $\Lambda(\cdot)$ is the logistic CDF, and \textit{elo\_fav} and \textit{elo\_und} are the Elo ratings of the stronger and the weaker team in the match, respectively.

\paragraph{Model 2: Logit for qualification change, group-level Elo ranks (Table~\ref{tab:reg_grp}).}
We replace the match-level Elo controls with the four group-level Elo ranks (\textit{elo1st} to \textit{elo4th}, from highest to lowest within the group), which describe the competitive balance of the whole group rather than of the specific match:
\[
\Pr(\textit{qual\_changed}_{gk} = 1)
  = \Lambda\!\left(\alpha + \beta\,\textit{FIFA\_rule}_g
  + \sum_{r=1}^{4} \gamma_r\,\textit{elo}^{(r)}_g\right)
\]

\paragraph{Model 3: Poisson for qualification count (Table~\ref{tab:reg_grp}).}
We collapse the data to one observation per group using the maximum value of \textit{qual\_count}, which is the cumulative number of times the set of qualifying teams changed within a group, and estimate a Poisson regression with robust standard errors:
\[
\mathbb{E}[\textit{qual\_count}_g]
  = \exp\!\left(\alpha + \beta\,\textit{FIFA\_rule}_g
  + \sum_{r=1}^{4} \gamma_r\,\textit{elo}^{(r)}_g\right)
\]

\paragraph{Models 4 and 5: Logit for suspense, goal level (Table~\ref{tab:reg_susp}).}
To test whether the FIFA rule is associated with higher suspense at the moment of each goal, we re-estimate the logit models with \textit{suspense} instead of \textit{qual\_changed} as the dependent variable, using the same goal-level dataset. Model~4 includes match-level Elo controls (\textit{elo\_fav}, \textit{elo\_und}), and Model~5 includes group-level Elo ranks (\textit{elo1st} to \textit{elo4th}).

\paragraph{MBM models: Logit for suspense and qualification change (Table~\ref{tab:reg_mbm}).}
At the minute-by-minute level, we estimate two logit models with group-level Elo controls (\textit{elo1st} to \textit{elo4th}). In MBM-A the dependent variable is the binary \textit{suspense} indicator. In MBM-B it is \textit{qual\_changed}, which tests whether the FIFA rule is associated with changes in the set of qualifying teams at the minute level.

Table~\ref{tab:reg_grp} reports the estimates of Models~1 to 3. The coefficient on \textit{FIFA\_rule} is negative in all specifications. World Cup goals are on average 9.7 percentage points less likely to change the set of qualifying teams than European Championship goals (Model~1, $p = 0.045$, $N = 607$). The sign is the same in Model~2 ($-6.5$ pp, $p = 0.146$, $N = 725$) and in the Poisson model for the qualification count (Model~3, coefficient $-0.33$, $p = 0.173$). MBM-B in Table~\ref{tab:reg_mbm} confirms this result at the minute level: World Cup minutes are 10.3 percentage points less likely to involve a change in the set of qualifying teams ($p = 0.034$), with similar results under the 3-points rule ($-0.105$, $p = 0.021$) and in the Elo common support sample ($-0.083$, $p = 0.082$).

The results are different when the dependent variable is \textit{suspense}, which indicates that one more goal \textit{would} change the set of qualifying teams. Table~\ref{tab:reg_susp} shows that the FIFA\_rule coefficient in Models~4 and 5 is small, changes sign, and is far from significant ($-0.015$, $p = 0.846$ in Model~4; $+0.019$, $p = 0.800$ in Model~5). The same holds for MBM-A ($-0.031$, $p = 0.696$). Thus, the FIFA rule reduces the number of \textit{actual} changes in the set of qualifying teams without changing the share of moments in which a change could occur. In other words, it affects realised turnover but not anticipated volatility.

The results differ between the 2-points and the 3-points periods. Under the old 2-points system (years $\le 1992$), the FIFA rule is associated with \textit{more} changes in the set of qualifying teams: $+0.290$ in MBM-B ($p = 0.041$), with positive but imprecise estimates in Models~1 to 3. Under the 3-points system (years $> 1992$), the sign becomes negative in all models ($-0.105$ in MBM-B, $p = 0.021$). One explanation is that draws were more common under the old system, so goal difference was more often decisive and groups under the FIFA rule saw more changes in the set of qualifying teams. When wins became more valuable and draws less common under the 3-points system, qualification was decided more often by points, which reduced the role of goal difference and the additional changes that the FIFA rule used to produce. Robustness checks are reported in the same tables.


\begin{table}[H]
\centering
\caption{Effect of FIFA tie-breaking rule on qualification change and count}
\label{tab:reg_grp}
\begin{threeparttable}
\small
\begin{tabular}{l cc cc cc}
\toprule
 & \multicolumn{2}{c}{M1: Logit (AME)} & \multicolumn{2}{c}{M2: Logit (AME)} & \multicolumn{2}{c}{M3: Poisson (coef)} \\
 & \multicolumn{2}{c}{\textit{qual\_changed}} & \multicolumn{2}{c}{\textit{qual\_changed}} & \multicolumn{2}{c}{\textit{qual\_count}} \\
\cmidrule(lr){2-3}\cmidrule(lr){4-5}\cmidrule(lr){6-7}
 & (1) & (2) & (3) & (4) & (5) & (6) \\
\midrule
\multicolumn{7}{l}{\textit{Panel A: Main specifications}} \\[2pt]
FIFA\_rule & $-0.097^{**}$ & $-0.083^{*}$ & $-0.065$ & $-0.043$ & $-0.331$ & $-0.271$ \\
          & $(0.048)$ & $(0.046)$ & $(0.045)$ & $(0.042)$ & $(0.243)$ & $(0.237)$ \\[3pt]
Match-level Elo & \checkmark & \checkmark & -- & -- & -- & -- \\
Group-level Elo & -- & -- & \checkmark & \checkmark & \checkmark & \checkmark \\
Year trend  & -- & \checkmark & -- & \checkmark & -- & \checkmark \\[3pt]
Observations    & 607 & 607 & 725 & 725 & 118 & 118 \\
\midrule
\multicolumn{7}{l}{\textit{Panel B: Robustness, FIFA\_rule only (all specifications include year trend)}} \\[2pt]
 & \multicolumn{2}{c}{M1} & \multicolumn{2}{c}{M2} & \multicolumn{2}{c}{M3} \\
\cmidrule(lr){2-3}\cmidrule(lr){4-5}\cmidrule(lr){6-7}
Two qualifying teams    & \multicolumn{2}{c}{$-0.074$} & \multicolumn{2}{c}{$-0.031$} & \multicolumn{2}{c}{$-0.118$} \\
                        & \multicolumn{2}{c}{$(0.046)$} & \multicolumn{2}{c}{$(0.044)$} & \multicolumn{2}{c}{$(0.304)$} \\
Three qualifying teams  & \multicolumn{2}{c}{$-0.414$} & \multicolumn{2}{c}{$-0.539$} & \multicolumn{2}{c}{$-2.597$} \\
                        & \multicolumn{2}{c}{$(0.629)$} & \multicolumn{2}{c}{$(0.499)$} & \multicolumn{2}{c}{$(1.892)$} \\
2-points rule ($\le 1992$) & \multicolumn{2}{c}{$0.044$} & \multicolumn{2}{c}{$0.242$} & \multicolumn{2}{c}{$0.599$} \\
                        & \multicolumn{2}{c}{$(0.125)$} & \multicolumn{2}{c}{$(0.170)$} & \multicolumn{2}{c}{$(0.829)$} \\
3-points rule ($> 1992$) & \multicolumn{2}{c}{$-0.092^{*}$} & \multicolumn{2}{c}{$-0.047$} & \multicolumn{2}{c}{$-0.284$} \\
                        & \multicolumn{2}{c}{$(0.048)$} & \multicolumn{2}{c}{$(0.041)$} & \multicolumn{2}{c}{$(0.252)$} \\
Elo common support      & \multicolumn{2}{c}{$-0.071$} & \multicolumn{2}{c}{$-0.043$} & \multicolumn{2}{c}{$-0.263$} \\
                        & \multicolumn{2}{c}{$(0.048)$} & \multicolumn{2}{c}{$(0.044)$} & \multicolumn{2}{c}{$(0.240)$} \\
\bottomrule
\end{tabular}
\begin{tablenotes}
\small
\item \textit{Notes:} M1 and M2 are logit models; coefficients are average marginal effects (AME). M3 is a Poisson model reporting log-rate coefficients. M1 controls for the Elo ratings of the stronger (\textit{elo\_fav}) and weaker (\textit{elo\_und}) team in each match. M2 and M3 include the four group-level Elo ranks (\textit{elo1st}--\textit{elo4th}). Standard errors are clustered at the group (year~$\times$~stage) level for M1--M2; HC0 robust standard errors for M3. The outcome variable is \textit{qual\_changed} (binary) in M1--M2 and the group-level maximum of \textit{qual\_count} in M3. Panel~B reports only the FIFA\_rule estimate; all Panel~B specifications include year trend. $^{*}p<0.10$, $^{**}p<0.05$, $^{***}p<0.01$.
\end{tablenotes}
\end{threeparttable}
\end{table}


\begin{table}[H]
\centering
\caption{Effect of FIFA tie-breaking rule on suspense, goal-level logit (AME)}
\label{tab:reg_susp}
\begin{threeparttable}
\small
\begin{tabular}{l cc cc}
\toprule
 & \multicolumn{2}{c}{M4: Match-level Elo} & \multicolumn{2}{c}{M5: Group-level Elo} \\
\cmidrule(lr){2-3}\cmidrule(lr){4-5}
 & (1) Main & (2) $+$Year & (3) Main & (4) $+$Year \\
\midrule
\multicolumn{5}{l}{\textit{Panel A: Main specifications}} \\[2pt]
FIFA\_rule & $-0.015$ & $-0.001$ & $\phantom{-}0.019$ & $\phantom{-}0.044$ \\
          & $(0.075)$ & $(0.077)$ & $(0.075)$ & $(0.076)$ \\[3pt]
Match-level Elo (\textit{elo\_fav}, \textit{elo\_und}) & \checkmark & \checkmark & -- & -- \\
Group-level Elo (\textit{elo1st}--\textit{elo4th})    & -- & -- & \checkmark & \checkmark \\
Year trend & -- & \checkmark & -- & \checkmark \\[3pt]
Observations & 607 & 607 & 725 & 725 \\
\midrule
\multicolumn{5}{l}{\textit{Panel B: Robustness, FIFA\_rule only (all specifications include year trend)}} \\[2pt]
 & \multicolumn{2}{c}{M4} & \multicolumn{2}{c}{M5} \\
\cmidrule(lr){2-3}\cmidrule(lr){4-5}
Two qualifying teams    & \multicolumn{2}{c}{$-0.080$} & \multicolumn{2}{c}{$-0.036$} \\
                        & \multicolumn{2}{c}{$(0.084)$} & \multicolumn{2}{c}{$(0.090)$} \\
Three qualifying teams  & \multicolumn{2}{c}{$-0.006$} & \multicolumn{2}{c}{$-0.734$} \\
                        & \multicolumn{2}{c}{$(0.642)$} & \multicolumn{2}{c}{$(0.630)$} \\
2-points rule ($\le 1992$) & \multicolumn{2}{c}{$-0.151$} & \multicolumn{2}{c}{$\phantom{-}0.079$} \\
                        & \multicolumn{2}{c}{$(0.160)$} & \multicolumn{2}{c}{$(0.314)$} \\
3-points rule ($> 1992$) & \multicolumn{2}{c}{$\phantom{-}0.031$} & \multicolumn{2}{c}{$\phantom{-}0.068$} \\
                        & \multicolumn{2}{c}{$(0.082)$} & \multicolumn{2}{c}{$(0.080)$} \\
Elo common support      & \multicolumn{2}{c}{$\phantom{-}0.006$} & \multicolumn{2}{c}{$\phantom{-}0.029$} \\
                        & \multicolumn{2}{c}{$(0.080)$} & \multicolumn{2}{c}{$(0.077)$} \\
\bottomrule
\end{tabular}
\begin{tablenotes}
\small
\item \textit{Notes:} Logit model for the binary \textit{suspense} indicator estimated on the goal-level dataset. Coefficients are average marginal effects (AME). M4 controls for match-level Elo ratings (\textit{elo\_fav}, \textit{elo\_und}); M5 uses group-level Elo ranks (\textit{elo1st}--\textit{elo4th}). Standard errors are clustered at the group (year~$\times$~stage) level. Panel~B robustness samples follow the same definitions as in Table~\ref{tab:reg_grp}; all Panel~B specifications include year trend. $^{*}p<0.10$, $^{**}p<0.05$, $^{***}p<0.01$.
\end{tablenotes}
\end{threeparttable}
\end{table}


\begin{table}[H]
\centering
\caption{MBM logit models: effect of FIFA tie-breaking rule on suspense and qualification change}
\label{tab:reg_mbm}
\begin{threeparttable}
\small
\begin{tabular}{l cc cc}
\toprule
 & \multicolumn{2}{c}{MBM-A: \textit{suspense}} & \multicolumn{2}{c}{MBM-B: \textit{qual\_changed}} \\
\cmidrule(lr){2-3}\cmidrule(lr){4-5}
 & (1) Main & (2) $+$Year & (3) Main & (4) $+$Year \\
\midrule
\multicolumn{5}{l}{\textit{Panel A: Main specifications}} \\[2pt]
FIFA\_rule & $-0.031$ & $\phantom{-}0.001$ & $-0.103^{**}$ & $-0.086^{*}$ \\
          & $(0.081)$ & $(0.083)$ & $(0.049)$ & $(0.045)$ \\[3pt]
Group-level Elo controls & \checkmark & \checkmark & \checkmark & \checkmark \\
Year trend & -- & \checkmark & -- & \checkmark \\[3pt]
Observations & 10{,}833 & 10{,}833 & 10{,}833 & 10{,}833 \\
\midrule
\multicolumn{5}{l}{\textit{Panel B: Robustness, FIFA\_rule only (all specifications include year trend)}} \\[2pt]
 & \multicolumn{2}{c}{MBM-A} & \multicolumn{2}{c}{MBM-B} \\
\cmidrule(lr){2-3}\cmidrule(lr){4-5}
Two qualifying teams    & \multicolumn{2}{c}{$-0.084$} & \multicolumn{2}{c}{$-0.049$} \\
                        & \multicolumn{2}{c}{$(0.095)$} & \multicolumn{2}{c}{$(0.052)$} \\
Three qualifying teams  & \multicolumn{2}{c}{$-0.219$} & \multicolumn{2}{c}{$-0.073$} \\
                        & \multicolumn{2}{c}{$(0.753)$} & \multicolumn{2}{c}{$(0.525)$} \\
2-points rule ($\le 1992$) & \multicolumn{2}{c}{$\phantom{-}0.265$} & \multicolumn{2}{c}{$\phantom{-}0.290^{**}$} \\
                        & \multicolumn{2}{c}{$(0.278)$} & \multicolumn{2}{c}{$(0.142)$} \\
3-points rule ($> 1992$) & \multicolumn{2}{c}{$\phantom{-}0.014$} & \multicolumn{2}{c}{$-0.105^{**}$} \\
                        & \multicolumn{2}{c}{$(0.086)$} & \multicolumn{2}{c}{$(0.046)$} \\
Elo common support      & \multicolumn{2}{c}{$-0.012$} & \multicolumn{2}{c}{$-0.083^{*}$} \\
                        & \multicolumn{2}{c}{$(0.084)$} & \multicolumn{2}{c}{$(0.048)$} \\
\bottomrule
\end{tabular}
\begin{tablenotes}
\small
\item \textit{Notes:} Logit models estimated on the minute-by-minute (MBM) dataset. MBM-A uses the binary \textit{suspense} indicator as the dependent variable; MBM-B uses \textit{qual\_changed}, a binary indicator equal to~1 when the qualifying composition changes at that minute. Group-level Elo ranks (\textit{elo1st}--\textit{elo4th}) are included as controls in all specifications. Standard errors are clustered at the group (year~$\times$~stage) level. Panel~B robustness samples follow the same definitions as in Table~\ref{tab:reg_grp}. $^{*}p<0.10$, $^{**}p<0.05$, $^{***}p<0.01$.
\end{tablenotes}
\end{threeparttable}
\end{table}



\section{Conclusion}
\label{sec:conclusion}

This paper studies how different tie-breaking criteria affect suspense in international football tournaments. Tie-breaking rules are not just technical details. They affect how attractive games are, with possible consequences for revenues, stadium attendance, and viewership. By choosing one criterion over another, governing bodies implicitly decide how much uncertainty remains during the group stage.

We compared the two main approaches in international football: the FIFA rule, which ranks goal difference first, and the UEFA rule, which ranks head-to-head results first. Our analysis combines two approaches. Paired tests within groups show that applying FIFA rules to the same group generates significantly more suspense than the UEFA criterion in the World Cup ($p < 0.01$), while the difference is not significant in the European Championship. A regression across the two tournaments that controls for competitive balance shows that World Cup goals are less likely to change the set of qualifying teams than European Championship goals. This suggests that the FIFA criterion spreads qualification uncertainty more evenly over the group stage, keeping more scenarios open at any given moment, instead of concentrating it in a few decisive goals.

The European Championship provides a useful contrast. The paired tests find no significant difference in suspense between the two rules. This suggests that the effect we find for the World Cup is specific to that competition and does not simply follow from how the two rules work in general.

Overall, our results show that institutional design choices, and in particular the choice of tie-breaking criterion, have measurable consequences for the suspense experienced by the very large audience that follows the World Cup. A rule that keeps more qualification scenarios open for longer increases the entertainment value of each match. Our evidence suggests that FIFA's goal-difference criterion generates more sustained group-level uncertainty. This is particularly relevant since FIFA has replaced goal difference with head-to-head results as the primary tie-breaker for the 2026 World Cup \citep{csato2026match}, and more generally for governing bodies as they continue to revise tournament formats to increase the commercial and sporting value of their competitions.


\bibliographystyle{apalike}
\bibliography{reference}


% ===========================================================================
\clearpage
\appendix
\section*{Appendix}
\renewcommand{\thesubsection}{A.\arabic{subsection}}
\setcounter{table}{0}
\setcounter{figure}{0}
\renewcommand{\thetable}{A\arabic{table}}
\renewcommand{\thefigure}{A\arabic{figure}}
% ===========================================================================

\subsection{Suspense Criteria: Full Logic}
\label{sec:suspense-criteria}

The following rules describe the exact criteria we use to assign a suspense value of 1 (suspenseful moment) in both FIFA and UEFA group stages. A suspenseful moment is defined as one where a single goal may change the composition of teams qualifying for the knockout stage. These rules are applied after each goal and at each match minute during the group stage.

\subsubsection{Suspense Criteria: FIFA}
\label{sec:suspense-fifa-implemented}

FIFA ranking within the group follows the tie-break hierarchy:
\[
\text{Points} \;\; \rightarrow \;\; \text{Overall Goal Difference} \;\; \rightarrow \;\; \text{Overall Goals Scored} \;\; \rightarrow \;\; \text{Head-to-Head (if required)}.
\]
We evaluate suspense after each goal/minute observation using the pipeline below.

\paragraph{Step 0: Identify the cut-off pair (leading vs lagging).}
Let \texttt{third\_qualify} indicate whether the relevant cut-off is:
\begin{itemize}
    \item \texttt{third\_qualify = 0}: \textbf{2nd} (leading, last qualified) vs \textbf{3rd} (lagging, first not qualified),
    \item \texttt{third\_qualify = 1}: \textbf{3rd} (leading) vs \textbf{4th} (lagging).
\end{itemize}
We denote the lagging and leading teams' current \emph{overall} differences by:
\[
\texttt{pts\_diff}_0,\;\texttt{gls\_diff}_0,\;\texttt{gs\_diff}_0,
\]
where \texttt{gs\_diff} is computed as goals scored by lagging minus goals scored by leading.

\paragraph{Step 1: Infer current match points from match goal difference.}
At every minute, each team has a match-specific goal difference (\texttt{last\_game\_goals\_diff}).
We infer match points robustly from this match GD:
\[
\texttt{match\_pts} =
\begin{cases}
\texttt{win\_pts} & \text{if match GD} > 0,\\
1 & \text{if match GD} = 0,\\
0 & \text{if match GD} < 0,
\end{cases}
\]
with \texttt{win\_pts} $=2$ if \texttt{year $\le$ 1992} and \texttt{win\_pts} $=3$ otherwise.

\paragraph{Step 2: Evaluate the two one-goal shocks }
% function calculate\_suspense.

We compute whether \emph{either} of the following shocks makes lagging rank $\ge$ leading under FIFA ordering:
\begin{enumerate}
    \item \textbf{Lagging scores (+1)}.
    \item \textbf{Leading concedes (+1)}.
\end{enumerate}

\noindent The shocks update \emph{both} (i) match states and (ii) overall differences:
\begin{itemize}
    \item If \texttt{h2h = 0} (teams are in different matches), the shock affects only the relevant team's match GD by $\pm 1$ and hence may change match points. Overall GD changes by $+1$ in favor of lagging for either shock. Goals scored changes by $+1$ \emph{only} if lagging scores.
    \item If \texttt{h2h = 1} (teams are facing each other in the current match), a single goal simultaneously improves lagging's match GD by $+1$ and worsens leading's match GD by $-1$. Hence overall GD changes by $+2$ in favor of lagging. Goals scored of lagging increases by $+1$.
\end{itemize}

\noindent After updating points (through inferred match points), overall GD, and overall goals scored, suspense is set to 1 if lagging becomes \emph{not behind} leading under:
\[
\text{Points} \rightarrow \text{Overall GD} \rightarrow \text{Overall GS}.
\]
If the shock produces a full tie on all three dimensions, we still classify the moment as suspenseful (\texttt{suspense=1}), since the next tie-break (e.g., head-to-head or lots) may decide qualification.

\paragraph{Step 3: Flag head-to-head tie-break cases }
% function update\_suspense
We create the indicator \texttt{h2h\_suspense} to identify the subset of observations where the qualification cut-off would be decided by the \emph{head-to-head} tie-break between the lagging and leading teams.
Practically, we flag \texttt{h2h\_suspense=1} when a full tie on the \emph{FIFA pre-H2H criteria} (points, overall goal difference, and overall goals scored) can be reached with a single-goal shock from \emph{both} perspectives:
\begin{itemize}
    \item full tie after \textbf{lagging scores (+1)}, and
    \item full tie after \textbf{leading concedes (+1)}.
\end{itemize}
Only in these cases do we set \texttt{suspense=1} and \texttt{h2h\_suspense=1}.
This ``both perspectives'' requirement excludes situations where only one shock produces a full tie, while the other shock would already give one team a strict advantage on goals scored.

\paragraph{Step 4: Evaluate the relevant head-to-head match}

% function \texttt{evaluate\_lagging\_losses\_and\_update\_suspense}.

For observations with \texttt{h2h\_suspense=1}, we evaluate the head-to-head status between the lagging and leading teams to determine whether the tie-break can actually favor the lagging team.
There are two distinct cases:
\begin{itemize}
    \item \textbf{\texttt{h2h = 1} (current match):} lagging and leading are facing each other in the ongoing match. In this case, the head-to-head outcome is determined by the \emph{current score} at minute $t$ (i.e., the current match goal difference for the lagging team).
    \item \textbf{\texttt{h2h = 0} (previous match):} lagging and leading are not playing each other in the current match minute; their head-to-head match is in one of the two \emph{previous} fixtures in the same group stage. In this case the head-to-head outcome is fixed and can be retrieved from \texttt{goals\_df} using the match's final score. If the lagging team loses the head-to-head match, we set \texttt{suspense} back to 0 for that observation.
\end{itemize}

\paragraph{Step 5: Draw-lots adjustments.}

% function \texttt{update\_suspense\_draw\_lots}
Finally, we manually set \texttt{suspense=1} in a small set of edge cases where the FIFA rules imply qualification may be decided by lots.

\paragraph{Diagnostic audit: outside-team threat when \texttt{suspense=0}}
% function (\texttt{audit\_suspense0\_where\_4th\_threatens\_2nd\_fifa}).
Because the main suspense logic evaluates only the cut-off pair (leading vs lagging), we also run a diagnostic check for cases with \texttt{suspense=0} and \texttt{third\_qualify=0} where the \textbf{4th} team can reach or overtake the \textbf{2nd} team with one goal shock under FIFA ordering.


\subsubsection{Suspense Criteria: UEFA}
\label{sec:suspense-uefa-implemented}

For competitions governed by UEFA rules, team ranking within a group follows a different tie-breaking hierarchy:
\begin{enumerate}
    \item total points,
    \item head-to-head result between the tied teams,
    \item overall goal difference,
    \item overall goals scored.
\end{enumerate}
Unlike the FIFA procedure, head-to-head is applied \emph{before} overall goal difference and goals scored.

\paragraph{Step 1: Identify the relevant ranking pair.}
At each observation (defined by a goal minute within a group-stage match), we identify the pair of teams defining the qualification cut-off. When only the top two teams qualify (\texttt{third\_qualify = 0}), the relevant comparison is between the teams ranked second and third. When the third-placed team may also qualify (\texttt{third\_qualify = 1}), the relevant comparison is between the teams ranked third and fourth.

\paragraph{Step 2: Evaluate one-goal shocks.}
For each observation we evaluate whether a single-goal shock could allow the lagging team to reach or overtake the leading team. Two possible shocks are considered: lagging scores one goal (+1) or leading concedes one goal (+1). Match points are inferred from the match goal difference using the appropriate points system.

\paragraph{Step 3: Apply UEFA tie-breaking rules.}
After applying the hypothetical shock, we compare the lagging and leading teams according to the UEFA ranking hierarchy: total points, then head-to-head, then overall goal difference, then overall goals scored. When the two teams are not currently facing each other, the head-to-head outcome is fixed from the previous fixture; when they are playing each other, the current score determines the head-to-head outcome.

\paragraph{Step 4: Determine suspense.}
If under either shock the lagging team becomes at least as well ranked as the leading team, we set \texttt{suspense = 1}.

\paragraph{Step 5: Identify outside-team threats.}
Finally, we check whether the fourth-placed team could reach or overtake the second-placed team with a single-goal shock, applying the full UEFA hierarchy. Observations satisfying this condition are edited separately.


\subsection{Correlation Analysis}
\label{app:correlations}

Tables~\ref{tab:corr_eu} to~\ref{tab:corr_wc_mbm} report pairwise correlations between suspense, the number of qualification changes, and group-level competitive balance (standard deviation of Elo ratings) for both tournaments and both datasets.

For the European Championship (Tables~\ref{tab:corr_eu} and~\ref{tab:corr_eu_mbm}), the number of changes in the set of qualifying teams is positively correlated with average suspense in every specification, and the correlation is stronger under FIFA rules. Suspense increases with group imbalance under both rules, while the correlation between imbalance and the number of qualification changes is weak and slightly negative. This is expected, since the qualification count captures only actual changes, whereas our suspense measure also captures potential changes.

\begin{table}[H]
\centering
\small
\caption{Correlation matrices for goal-level aggregation in the UEFA European Championship}
\label{tab:corr_eu}
\begin{tabular}{lccc}
\toprule
 & Elo std & Avg. qual. count & Avg. suspense \\
\midrule
\multicolumn{4}{l}{\textbf{UEFA}} \\
Elo std             & 1.000 & -0.089 & 0.251 \\
Avg. qual. count    & -0.089 & 1.000 & 0.391 \\
Avg. suspense       & 0.251 & 0.391 & 1.000 \\
\midrule
\multicolumn{4}{l}{\textbf{FIFA}} \\
Elo std             & 1.000 & -0.061 & 0.264 \\
Avg. qual. count    & -0.061 & 1.000 & 0.406 \\
Avg. suspense       & 0.264 & 0.406 & 1.000 \\
\bottomrule
\end{tabular}
\end{table}

\begin{table}[H]
\centering
\small
\caption{Correlation matrices for minute-by-minute aggregation in the UEFA European Championship}
\label{tab:corr_eu_mbm}
\begin{tabular}{lccc}
\toprule
 & Elo std & Avg. qual. count & Avg. suspense \\
\midrule
\multicolumn{4}{l}{\textbf{UEFA}} \\
Elo std             & 1.000 & -0.102 & 0.175 \\
Avg. qual. count    & -0.102 & 1.000 & 0.299 \\
Avg. suspense       & 0.175 & 0.299 & 1.000 \\
\midrule
\multicolumn{4}{l}{\textbf{FIFA}} \\
Elo std             & 1.000 & -0.036 & 0.153 \\
Avg. qual. count    & -0.036 & 1.000 & 0.351 \\
Avg. suspense       & 0.153 & 0.351 & 1.000 \\
\bottomrule
\end{tabular}
\end{table}


For the World Cup (Tables~\ref{tab:corr_wc} and~\ref{tab:corr_wc_mbm}), the number of changes in the set of qualifying teams is positively correlated with average suspense (correlations between 0.40 and 0.59), and slightly more so under FIFA rules. Unlike in the European Championship, suspense is negatively correlated with group imbalance, which means that more balanced groups tend to generate more suspense.

\begin{table}[H]
\centering
\small
\caption{Correlation matrices for goal-level aggregation in the FIFA World Cup}
\label{tab:corr_wc}
\begin{tabular}{lccc}
\toprule
 & Elo std & Avg. qual. count & Avg. suspense \\
\midrule
\multicolumn{4}{l}{\textbf{UEFA rule}} \\
Elo std             & 1.000 & -0.064 & -0.184 \\
Avg. qual. count    & -0.064 & 1.000 & 0.456 \\
Avg. suspense       & -0.184 & 0.456 & 1.000 \\
\midrule
\multicolumn{4}{l}{\textbf{FIFA rule}} \\
Elo std             & 1.000 & -0.060 & -0.278 \\
Avg. qual. count    & -0.060 & 1.000 & 0.588 \\
Avg. suspense       & -0.278 & 0.588 & 1.000 \\
\bottomrule
\end{tabular}
\end{table}

\begin{table}[H]
\centering
\small
\caption{Correlation matrices for minute-by-minute aggregation in the FIFA World Cup}
\label{tab:corr_wc_mbm}
\begin{tabular}{lccc}
\toprule
 & Elo std & Avg. qual. count & Avg. suspense \\
\midrule
\multicolumn{4}{l}{\textbf{UEFA rule}} \\
Elo std             & 1.000 & -0.006 & -0.222 \\
Avg. qual. count    & -0.006 & 1.000 & 0.403 \\
Avg. suspense       & -0.222 & 0.403 & 1.000 \\
\midrule
\multicolumn{4}{l}{\textbf{FIFA rule}} \\
Elo std             & 1.000 & -0.006 & -0.276 \\
Avg. qual. count    & -0.006 & 1.000 & 0.476 \\
Avg. suspense       & -0.276 & 0.476 & 1.000 \\
\bottomrule
\end{tabular}
\end{table}




\subsection{World Cup Suspense Trajectories by Group Type}
\label{app:wc_trajectories}

To examine whether the FIFA advantage leads to suspense that is spread more evenly over the match, we classify World Cup groups as \textit{stable} or \textit{volatile}, depending on whether their average number of qualification changes is below or above the median, and plot the minute-by-minute suspense trajectory separately for each type.

Volatile groups have higher suspense than stable groups under both rules. Under FIFA rules (Figure~\ref{fig:susp_traj_wc_fifa}), the gap between stable and volatile groups is already large at kick-off (around 0.4) and remains large throughout the match, as both trajectories decline gradually. Under UEFA rules (Figure~\ref{fig:susp_traj_wc_uefa}), the initial gap is smaller (about 0.2) but widens considerably as the match goes on.

\begin{figure}[H]
    \centering
    \includegraphics[width=1\textwidth]{images/susp_traj_wc_fifa.png}
    \caption{Suspense trajectories in World Cup groups (stable vs.\ volatile) under FIFA rules.}
    \label{fig:susp_traj_wc_fifa}
\end{figure}

\begin{figure}[H]
    \centering
    \includegraphics[width=1\textwidth]{images/susp_traj_wc_uefa.png}
    \caption{Suspense trajectories in World Cup groups (stable vs.\ volatile) under UEFA rules.}
    \label{fig:susp_traj_wc_uefa}
\end{figure}


\subsection{European Championship Suspense Trajectories}
\label{app:euro_trajectories}

Figures~\ref{fig:susp_traj_eu_uefa} and~\ref{fig:susp_traj_eu_fifa} show the suspense trajectories of stable and volatile groups in the European Championship under UEFA and FIFA rules, respectively. Under UEFA rules (Figure~\ref{fig:susp_traj_eu_uefa}), the difference between the two trajectories is small at the start and at the end of matches and largest between the 30th and 50th minute. Under FIFA rules (Figure~\ref{fig:susp_traj_eu_fifa}), the gap starts at about 0.3 and narrows gradually toward 0.2. In both cases, the differences are smaller than in the World Cup, which is consistent with the non-significant results of the statistical tests and the regressions.

\begin{figure}[H]
    \centering
    \includegraphics[width=1\textwidth]{images/susp_traj_eu_uefa.png}
    \caption{Suspense trajectories in European Championship groups (stable vs.\ volatile) under UEFA rules.}
    \label{fig:susp_traj_eu_uefa}
\end{figure}

\begin{figure}[H]
    \centering
    \includegraphics[width=1\textwidth]{images/susp_traj_eu_fifa.png}
    \caption{Suspense trajectories in European Championship groups (stable vs.\ volatile) under FIFA rules.}
    \label{fig:susp_traj_eu_fifa}
\end{figure}


\end{document}
